\documentclass[10pt,a4paper]{amsart}
\usepackage[margin=19mm]{geometry}
\usepackage{amsmath,amssymb,mathtools}
\usepackage[scr=rsfs,cal=euler]{mathalfa}
\usepackage{xfrac}
\usepackage[unicode,hidelinks,bookmarks=false]{hyperref}
\newcommand{\ie}{i.e.\ }
\newcommand{\cf}{cf.\ }
\newcommand{\resp}{respectively\ }
\newcommand{\Subsection}{\subsection}
\providecommand{\nobreakdash}{\nobreak\hbox{-}}
\setlength{\emergencystretch}{2em}
\multlinegap0pt
\usepackage{bm}
\usepackage{bbm}
\usepackage{dsfont}
\usepackage{xspace}
\usepackage{cancel}
\usepackage{mathtools}
\usepackage{amsthm}
\makeatletter
\@ifundefined{lem}{\newtheorem{lem}{Lemma}}{}
\makeatother
\makeatletter
\expandafter\ifx\csname mythm\endcsname\relax
\newenvironment{mythm}[1]
  {\renewcommand\theinnercustomthm{#1}\innercustomthm}
  {\endinnercustomthm}
\fi
\makeatother
\title[Proof complexity of Mal'tsev CSP]{Proof complexity of Mal'tsev CSP}
\author{Azza Gaysin}
\date{Source: arXiv:2508.00396v1 (CC BY 4.0)}
\begin{document}
\maketitle

\noindent\emph{Source and licence.} Proof complexity of Mal'tsev CSP. Version arXiv:2508.00396v1 (CC BY 4.0), distributed under the Creative Commons Attribution 4.0 International licence, as recorded in the arXiv licence metadata for this exact version and shown on the abstract page https://arxiv.org/abs/2508.00396v1. Full rights notice: licenses/arxiv-2508.00396-CC-BY-4.0.txt.\par\smallskip

\noindent\textit{Editorial scope.} Counted unit: the Lem whose source label is \texttt{alksjdca;syfrg} in the article (source0.tex lines 1098--1100, source label \texttt{alksjdca;syfrg}), delivered together with the article's own complete original proof of it (source0.tex lines 1102--1119, one \texttt{proof} environment of 18 lines). No mathematics is written, repaired or re-derived: statement and proof are the article's own text line for line. Required context retained uncounted: simplemm (lines 807--809); =a=-sify34gtth (lines 1089--1096). Every definition, display and label of those lines is retained unchanged. Omitted from this package and not redistributed: 1--806, 810--1088, 1097--1097, 1101--1101, 1120--1450. Nothing inside a retained excerpt is omitted or altered: every retained body is delivered exactly as the article's lines are written, so no omission receipt of any kind accompanies this package. Citations: the retained excerpts carry no citation command at all, so no bibliography entry of the article is transcribed and no reference is redistributed. Reference adaptation: none was needed and none was made; every reference inside the delivered unit resolves inside it, so no unresolved reference or citation is produced. Printed slips kept literal: no printed word, symbol, hypothesis or number of the counted statement or of its proof was altered. Layout and class: the article's own class is \texttt{article} with packages including geometry, algcompatible, algorithm, algpseudocode, inputenc, fontenc, lmodern, amsmath, amssymb, amsfonts, amsthm, float, hyperref, rsfso; the wrapper is the lane's pinned \texttt{amsart} editorial wrapper at 10pt on A4, which restates the theorem-like environments the retained text needs and adds the article's own macro definitions that the retained text needs (none), with extra wrapper packages (bm, bbm, dsfont, xspace, cancel, mathtools). The delivered numbering is the unit's own. Assets: no retained span contains a figure environment, a graphics-inclusion command, a PDF-page inclusion, a table or a TikZ picture; no image file is copied into this package, the package declares no asset dependency and no image file is redistributed with it. Licence: version arXiv:2508.00396v1 of this article is distributed under the Creative Commons Attribution 4.0 International licence, as recorded in the arXiv licence metadata for this exact version and shown on the abstract page https://arxiv.org/abs/2508.00396v1. A full rights notice is delivered at licenses/arxiv-2508.00396-CC-BY-4.0.txt. Characters: every retained excerpt is pure ASCII, so the delivered engine, XeLaTeX with its default Latin Modern text font, reports no missing character.
\begin{lem}\label{simplemm}
$V^1$ proves that for any CSP instance $\Theta$ corresponding to algebra $([q], M)$, and any homomorphisms $H_1, H_2, H_3$ from $\mathcal{X}$ to $\ddot{\mathcal{A}}$, the map $H = \mu(H_1, H_2, H_3)$ is again a homomorphism from $\mathcal{X}$ to $\ddot{\mathcal{A}}$.
\end{lem}

\begin{lem}[Correctness of $nonempty^3$]\label{=a=-sify34gtth}
Assume that $U^{[j-1]} \subseteq \mathscr{U}_{\Theta'_{j-1}}$ and $\mathrm{Sig}(U^{[j-1]}) = \mathrm{Sig}(\mathscr{U}_{\Theta'_{j-1}})$. Then 
$V^1$ proves that the function $nonempty^3(U^{[j-1]}, j, i, \{(a_j, a)\})$ for any $i > j$, any $a_j \in D_j$, and any $a \in D_i$ either outputs a homomorphism $T$ from $\mathcal{X}'_{j-1}$ to $\ddot{\mathcal{A}}'$ such that $\pi_{j,i}T = (a_j, a)$, or, if
\[
nonempty^3(U^{[j-1]}, j, i, \{(a_j, a)\}) = \emptyset,
\]
then there is no homomorphism $T \in \mathscr{U}_{\Theta'_{j-1}}$ such that $\pi_{j,i}T = (a_j, a)$.
\end{lem}

\begin{lem}[Correctness of $fixvalues$]\label{alksjdca;syfrg}
Assume that $R^{[l-1]} \subseteq \mathscr{R}_{\Theta_{l-1}}$ and $\mathrm{Sig}(R^{[l-1]}) = \mathrm{Sig}(\mathscr{R}_{\Theta_{l-1}})$, $H_a\in \mathscr{R}_{\Theta_{l-1}}$, and $k < n$. Then $V^1$ proves that the function $fixvalues(R^{[l-1]}, H_a, k-1)$ outputs a compact representation $U^{[k-1]}$ of the instance $\mathscr{U}_{\Theta'_{k-1}}$, i.e., $U^{[k-1]} \subseteq \mathscr{U}_{\Theta'_{k-1}}$ and $\mathrm{Sig}(U^{[k-1]}) = \mathrm{Sig}(\mathscr{U}_{\Theta'_{k-1}})$.
\end{lem}

\begin{proof}
First, by construction, we have $U^{[-1]} \subseteq \mathscr{U}_{\Theta'_{-1}}$ and $\mathrm{Sig}(U^{[-1]}) = \mathrm{Sig}(\mathscr{U}_{\Theta'_{-1}})$. To prove the lemma, it suffices to show that for every $0 \leq j < k-1$, if 
\[
U^{[j-1]} \subseteq \mathscr{U}_{\Theta'_{j-1}} \quad \text{and} \quad \mathrm{Sig}(U^{[j-1]}) = \mathrm{Sig}(\mathscr{U}_{\Theta'_{j-1}}),
\]
then
\[
U^{[j]} \subseteq \mathscr{U}_{\Theta'_{j}} \quad \text{and} \quad \mathrm{Sig}(U^{[j]}) = \mathrm{Sig}(\mathscr{U}_{\Theta'_{j}}).
\]

By the construction of the instances, it is clear that $\mathscr{U}_{\Theta'_{j}} \subseteq \mathscr{U}_{\Theta'_{j-1}}$, and hence $\mathrm{Sig}(\mathscr{U}_{\Theta'_{j}}) \subseteq \mathrm{Sig}(\mathscr{U}_{\Theta'_{j-1}})$. The only maps that can belong to $\mathscr{U}_{\Theta'_{j}}$ are those that send $0$ to $a_0$, \ldots , and $j$ to $a_j$. Therefore, if $\mathscr{U}_{\Theta'_{j}}$ is non-empty, the tuple $(j, a_j, a_j)$ must belong to $\mathrm{Sig}(U^{[j-1]})$. We copy to $U^{[j]}$ a map $T$ with this index only if such a map exists in $U^{[j-1]}$. Since $U^{[j-1]} \subseteq \mathscr{U}_{\Theta'_{j-1}}$ and $T(j) = a_j$, it follows that $T \in \mathscr{U}_{\Theta'_{j}}$. This is the only map we copy to $U^{[j]}$ from $U^{[j-1]}$ for $i \leq j$. Note that no other maps are needed for these coordinates, as we do not need to generate additional values.

For all $(i, a, b) \in \mathrm{Sig}(\mathscr{U}_{\Theta'_{j-1}})$ with $i > j$, we store in $U^{[j i a b]}$ the map $$nonempty^3(U^{[j-1]}, j, i, \{(a_j, a)\}),$$ and in $U^{[j i b a]}$ the map 
\[
\mu(nonempty^3(U^{[j-1]}, j, i, \{(a_j, a)\}), U^{[(j-1) i a b]}, U^{[(j-1) i b a]}).
\]
In the first case, by Lemma~\ref{=a=-sify34gtth}, this is a homomorphism $T$ from $\mathscr{U}_{\Theta'_{j-1}}$ such that $\pi_{j i} T = (a_j, a)$, and hence $T \in \mathscr{U}_{\Theta'_{j}}$. In the second case, by Lemma~\ref{simplemm}, the result is a homomorphism $T$ from $\mathscr{U}_{\Theta'_{j}}$ such that $\pi_{j i} T = (a_j, b)$. Thus, $U^{[j]} \subseteq \mathscr{U}_{\Theta'_{j}}$. The equality $\mathrm{Sig}(U^{[j]}) = \mathrm{Sig}(\mathscr{U}_{\Theta'_{j}})$ follows from the correctness of the function $nonempty^3$.
\end{proof}

\end{document}
